% Algebra 2 Lesson 6-6 fast-track packet: Exponential and Logarithmic Equations
% (enVision A2 2024 TE, Example 1 and Example 5 only; proposal, PLC has not met on this lesson)
% Shared preamble for the "fast" (hacked-down) Algebra 2 lesson packets, Oct 2026.
% Each packet: % Shared preamble for the "fast" (hacked-down) Algebra 2 lesson packets, Oct 2026.
% Each packet: % Shared preamble for the "fast" (hacked-down) Algebra 2 lesson packets, Oct 2026.
% Each packet: \input{fast-packet-template} then \begin{document} ... \end{document}.
\documentclass[12pt]{article}
\usepackage[margin=0.7in]{geometry}
\usepackage{amsmath,amssymb}
\usepackage{booktabs,tabularx,array}
\usepackage{enumitem}
\usepackage[T1]{fontenc}
\usepackage[utf8]{inputenc}
\usepackage{lmodern}
\usepackage{tikz}
\usepackage{pgfplots}
\pgfplotsset{compat=1.18}
\usepackage{xcolor}
\setlist{nosep,leftmargin=1.4em}
\setlength{\parindent}{0pt}
\setlength{\parskip}{6pt}
\renewcommand{\arraystretch}{1.35}
\newcolumntype{Y}{>{\raggedright\arraybackslash}X}
\pagestyle{empty}

% blanks and workspace
\newcommand{\blank}[1][2.6cm]{\rule[-2pt]{#1}{0.5pt}}
\newcommand{\lines}[1][1]{\par\vspace{2pt}\foreach \i in {1,...,#1}{\rule{0pt}{22pt}\rule{\linewidth}{0.4pt}\par}}
\newcommand{\wspart}[2]{\vspace{4pt}{\large\bfseries #1: #2}\par\vspace{-6pt}\rule{\textwidth}{0.4pt}\par}

% header: \packethead{Lesson 1-3}{Piecewise-Defined Functions}
\newcommand{\packethead}[2]{%
{\LARGE\bfseries #2}\hfill Name \blank[5cm]\\[2pt]
Algebra 2 \quad #1 (fast track) \hfill Date \blank[2.5cm] \quad Period \blank[1.2cm]\par\vspace{2pt}}

% boxed goals: \goals{learning goal}{language goal}
\newcommand{\goals}[2]{%
\fbox{\parbox{\dimexpr\textwidth-2\fboxsep-2\fboxrule}{%
\textbf{Learning goal:} #1\\[2pt]
\textbf{Language goal:} #2}}\par}

% quiz target box: \quiztarget{text}
\newcommand{\quiztarget}[1]{%
\vspace{4pt}\colorbox{gray!12}{\parbox{\dimexpr\textwidth-2\fboxsep}{\textbf{On the quiz:} #1}}\par}

% IXL table row helper: code & skill & which quiz question it feeds
\newcommand{\ixlrow}[3]{\texttt{#1} & #2 & #3 \\}

% one grid style for every graph; override xmin/xmax/... per plot
\pgfplotsset{sheetgrid/.style={
  axis lines=middle, grid=both, grid style={gray!45},
  xmin=-10.5, xmax=10.5, ymin=-10.5, ymax=10.5,
  xtick={-10,-8,...,10}, ytick={-10,-8,...,10},
  x=0.28cm, y=0.28cm,
  ticklabel style={font=\scriptsize, fill=white, inner sep=1pt},
  xlabel={$x$}, ylabel={$y$},
  axis line style={thick}, clip=true}}
% unit grid for small windows
\pgfplotsset{unitgrid/.style={sheetgrid,
  xmin=-5.5, xmax=5.5, ymin=-5.5, ymax=5.5,
  xtick={-5,...,5}, ytick={-5,...,5}, x=0.5cm, y=0.5cm}}
% open and closed endpoints
\tikzset{closed/.style={circle, fill, inner sep=1.6pt},
         open/.style={circle, draw, fill=white, inner sep=1.4pt}}
 then \begin{document} ... \end{document}.
\documentclass[12pt]{article}
\usepackage[margin=0.7in]{geometry}
\usepackage{amsmath,amssymb}
\usepackage{booktabs,tabularx,array}
\usepackage{enumitem}
\usepackage[T1]{fontenc}
\usepackage[utf8]{inputenc}
\usepackage{lmodern}
\usepackage{tikz}
\usepackage{pgfplots}
\pgfplotsset{compat=1.18}
\usepackage{xcolor}
\setlist{nosep,leftmargin=1.4em}
\setlength{\parindent}{0pt}
\setlength{\parskip}{6pt}
\renewcommand{\arraystretch}{1.35}
\newcolumntype{Y}{>{\raggedright\arraybackslash}X}
\pagestyle{empty}

% blanks and workspace
\newcommand{\blank}[1][2.6cm]{\rule[-2pt]{#1}{0.5pt}}
\newcommand{\lines}[1][1]{\par\vspace{2pt}\foreach \i in {1,...,#1}{\rule{0pt}{22pt}\rule{\linewidth}{0.4pt}\par}}
\newcommand{\wspart}[2]{\vspace{4pt}{\large\bfseries #1: #2}\par\vspace{-6pt}\rule{\textwidth}{0.4pt}\par}

% header: \packethead{Lesson 1-3}{Piecewise-Defined Functions}
\newcommand{\packethead}[2]{%
{\LARGE\bfseries #2}\hfill Name \blank[5cm]\\[2pt]
Algebra 2 \quad #1 (fast track) \hfill Date \blank[2.5cm] \quad Period \blank[1.2cm]\par\vspace{2pt}}

% boxed goals: \goals{learning goal}{language goal}
\newcommand{\goals}[2]{%
\fbox{\parbox{\dimexpr\textwidth-2\fboxsep-2\fboxrule}{%
\textbf{Learning goal:} #1\\[2pt]
\textbf{Language goal:} #2}}\par}

% quiz target box: \quiztarget{text}
\newcommand{\quiztarget}[1]{%
\vspace{4pt}\colorbox{gray!12}{\parbox{\dimexpr\textwidth-2\fboxsep}{\textbf{On the quiz:} #1}}\par}

% IXL table row helper: code & skill & which quiz question it feeds
\newcommand{\ixlrow}[3]{\texttt{#1} & #2 & #3 \\}

% one grid style for every graph; override xmin/xmax/... per plot
\pgfplotsset{sheetgrid/.style={
  axis lines=middle, grid=both, grid style={gray!45},
  xmin=-10.5, xmax=10.5, ymin=-10.5, ymax=10.5,
  xtick={-10,-8,...,10}, ytick={-10,-8,...,10},
  x=0.28cm, y=0.28cm,
  ticklabel style={font=\scriptsize, fill=white, inner sep=1pt},
  xlabel={$x$}, ylabel={$y$},
  axis line style={thick}, clip=true}}
% unit grid for small windows
\pgfplotsset{unitgrid/.style={sheetgrid,
  xmin=-5.5, xmax=5.5, ymin=-5.5, ymax=5.5,
  xtick={-5,...,5}, ytick={-5,...,5}, x=0.5cm, y=0.5cm}}
% open and closed endpoints
\tikzset{closed/.style={circle, fill, inner sep=1.6pt},
         open/.style={circle, draw, fill=white, inner sep=1.4pt}}
 then \begin{document} ... \end{document}.
\documentclass[12pt]{article}
\usepackage[margin=0.7in]{geometry}
\usepackage{amsmath,amssymb}
\usepackage{booktabs,tabularx,array}
\usepackage{enumitem}
\usepackage[T1]{fontenc}
\usepackage[utf8]{inputenc}
\usepackage{lmodern}
\usepackage{tikz}
\usepackage{pgfplots}
\pgfplotsset{compat=1.18}
\usepackage{xcolor}
\setlist{nosep,leftmargin=1.4em}
\setlength{\parindent}{0pt}
\setlength{\parskip}{6pt}
\renewcommand{\arraystretch}{1.35}
\newcolumntype{Y}{>{\raggedright\arraybackslash}X}
\pagestyle{empty}

% blanks and workspace
\newcommand{\blank}[1][2.6cm]{\rule[-2pt]{#1}{0.5pt}}
\newcommand{\lines}[1][1]{\par\vspace{2pt}\foreach \i in {1,...,#1}{\rule{0pt}{22pt}\rule{\linewidth}{0.4pt}\par}}
\newcommand{\wspart}[2]{\vspace{4pt}{\large\bfseries #1: #2}\par\vspace{-6pt}\rule{\textwidth}{0.4pt}\par}

% header: \packethead{Lesson 1-3}{Piecewise-Defined Functions}
\newcommand{\packethead}[2]{%
{\LARGE\bfseries #2}\hfill Name \blank[5cm]\\[2pt]
Algebra 2 \quad #1 (fast track) \hfill Date \blank[2.5cm] \quad Period \blank[1.2cm]\par\vspace{2pt}}

% boxed goals: \goals{learning goal}{language goal}
\newcommand{\goals}[2]{%
\fbox{\parbox{\dimexpr\textwidth-2\fboxsep-2\fboxrule}{%
\textbf{Learning goal:} #1\\[2pt]
\textbf{Language goal:} #2}}\par}

% quiz target box: \quiztarget{text}
\newcommand{\quiztarget}[1]{%
\vspace{4pt}\colorbox{gray!12}{\parbox{\dimexpr\textwidth-2\fboxsep}{\textbf{On the quiz:} #1}}\par}

% IXL table row helper: code & skill & which quiz question it feeds
\newcommand{\ixlrow}[3]{\texttt{#1} & #2 & #3 \\}

% one grid style for every graph; override xmin/xmax/... per plot
\pgfplotsset{sheetgrid/.style={
  axis lines=middle, grid=both, grid style={gray!45},
  xmin=-10.5, xmax=10.5, ymin=-10.5, ymax=10.5,
  xtick={-10,-8,...,10}, ytick={-10,-8,...,10},
  x=0.28cm, y=0.28cm,
  ticklabel style={font=\scriptsize, fill=white, inner sep=1pt},
  xlabel={$x$}, ylabel={$y$},
  axis line style={thick}, clip=true}}
% unit grid for small windows
\pgfplotsset{unitgrid/.style={sheetgrid,
  xmin=-5.5, xmax=5.5, ymin=-5.5, ymax=5.5,
  xtick={-5,...,5}, ytick={-5,...,5}, x=0.5cm, y=0.5cm}}
% open and closed endpoints
\tikzset{closed/.style={circle, fill, inner sep=1.6pt},
         open/.style={circle, draw, fill=white, inner sep=1.4pt}}

\begin{document}

{\let\LARGE\large\packethead{Lesson 6-6}{Exponential and Logarithmic Equations}}
\goals{I can solve an exponential equation by rewriting both sides with a common base, and solve a logarithmic equation by setting the arguments equal and checking for extraneous solutions.}{I can explain each step of my solution using the words \emph{common base}, \emph{exponent}, \emph{argument} and \emph{extraneous solution}.}

{\small Spanish: exponential equation = \emph{ecuaci\'on exponencial}; logarithmic equation = \emph{ecuaci\'on logar\'itmica}.}

\quiztarget{You will (1) solve an exponential equation with no common base by taking the log of both sides (full Savvas quiz only; optional IXL 9F2); (2) rewrite two clouds' exponential expressions, which have different bases, with a common base to find when they are equal, then find the volume (Example 1; the bonus); (3) find and fix the mistake in someone's common-base work and name each step (Example 1 and the Practice); (4) solve a log equation with a coefficient in front of one log (move it inside first) and keep only the solutions that make every argument positive (Example 5); (5) solve a log equation by graphing (full Savvas quiz only).}

\wspart{Example 1}{Solve Exponential Equations Using a Common Base}
\textbf{Property of Equality for Exponential Equations:} if $b>0$ and $b\neq1$, then $b^x=b^y$ if and only if $x=y$. Two powers of the same base are equal exactly when their \textbf{exponents} are equal.

What is the solution to $\left(\tfrac12\right)^{x+7}=4^{3x}$? \quad The key: $\tfrac12$ and $4$ are both powers of 2, so 2 is a \textbf{common base}.

\begin{center}
\begin{tabular}{|l|l|l|}
\hline
\textbf{Line} & \textbf{Work} & \textbf{Why} \\ \hline
1 & $\left(\tfrac12\right)^{x+7}=4^{3x}$ & the equation \\
2 & $\left(2^{-1}\right)^{x+7}=\left(2^{2}\right)^{3x}$ & rewrite each side with a common base: $\tfrac12=2^{-1}$, $4=2^2$ \\
3 & $2^{-x-7}=2^{6x}$ & Power of a Power Property: multiply the exponents \\
4 & $-x-7=6x$ & Property of Equality for Exponential Equations \\
5 & $-7=7x$ & add $x$ to both sides \\
6 & $x=-1$ & divide by 7 \\ \hline
\end{tabular}
\end{center}

\textbf{Check} $x=-1$: \ $\left(\tfrac12\right)^{-1+7}=\left(\tfrac12\right)^{6}=\tfrac1{64}$ \ and \ $4^{3(-1)}=4^{-3}=\tfrac1{64}$. \checkmark

\textbf{Watch line 3.} $\left(2^{-1}\right)^{x+7}$ multiplies $-1$ by the \emph{whole} exponent $x+7$: $-x-7$, not $-x+7$.

\textbf{Say it:} I wrote both sides with the common base \blank[1.5cm]. Then I set the exponents \blank[2.2cm] equal because \dots
\lines[2]

\newpage

\wspart{Try It 1}{Solve each equation using a common base}
Hint: for each part, name the common base, rewrite both sides, set the exponents equal, solve, and check.

\begin{minipage}[t]{0.48\textwidth}
\textbf{a.} $25^{3x}=125^{x+2}$

\vspace{2pt}
Common base: \blank[2cm]
\lines[6]
$x=\blank[2cm]$ \quad Check: \blank[3.2cm]
\end{minipage}\hfill
\begin{minipage}[t]{0.48\textwidth}
\textbf{b.} $0.001=10^{6x}$

\vspace{2pt}
Common base: \blank[2cm]
\lines[6]
$x=\blank[2cm]$ \quad Check: \blank[3.2cm]
\end{minipage}

% R5 Practice: structural parallel of Lesson Quiz item 3 (Max, 5^{x+3}=25^{x-2}, error in the
% Power of a Power line). Different base and numbers; never the quiz item itself.
\wspart{Practice}{Find the mistake}
Lena solved $2^{x+5}=8^{x-1}$, but she made an error. Her work is shown.

\begin{minipage}[t]{0.52\textwidth}
\vspace{0pt}
\begin{tabular}{|c|l|l|}
\hline
\textbf{Line} & \textbf{Lena's work} & \textbf{Step} \\ \hline
1 & $8^{x-1}=2^{x+5}$ & \\[6pt]
2 & $\left(2^{3}\right)^{x-1}=2^{x+5}$ & \blank[2.6cm] \\[6pt]
3 & $2^{3x-1}=2^{x+5}$ & \blank[2.6cm] \\[6pt]
4 & $3x-1=x+5$ & \blank[2.6cm] \\[6pt]
5 & $x=3$ & \\ \hline
\end{tabular}
\end{minipage}\hfill
\begin{minipage}[t]{0.45\textwidth}
\vspace{0pt}\raggedright
\textbf{a.} Name the step in lines 2, 3 and 4. Choose from:
\begin{itemize}
\item Rewrite with a common base
\item Power of a Power Property
\item Property of Equality for Exponential Equations
\end{itemize}
\end{minipage}

\vspace{4pt}
\textbf{b.} Lena made an error in line \blank[1.2cm]. The line should be \blank[4cm]. The correct solution is $x=\blank[1.5cm]$.

\textbf{c.} Check your solution in $2^{x+5}=8^{x-1}$:
\lines[3]

\newpage

\wspart{Example 5}{Solve Logarithmic Equations}
\textbf{Property of Equality for Logarithmic Equations:} if $b>0$, $b\neq1$, and $x$ and $y$ are positive, then $\log_b x=\log_b y$ if and only if $x=y$: the \textbf{arguments} (the expressions inside the logs) are equal.

What is the solution to $\ln(x^2-16)=\ln(6x)$?

\begin{minipage}[t]{0.52\textwidth}
\vspace{0pt}
\begin{tabular}{|l|>{\raggedright\arraybackslash}p{4.75cm}|}
\hline
\textbf{Work} & \textbf{Why} \\ \hline
$\ln(x^2-16)=\ln(6x)$ & the equation \\
$x^2-16=6x$ & Property of Equality for Logarithmic Equations \\
$x^2-6x-16=0$ & subtract $6x$ \\
$(x-8)(x+2)=0$ & factor \\
$x=8$ \ or \ $x=-2$ & Zero Product Property \\ \hline
\end{tabular}
\end{minipage}\hfill
\begin{minipage}[t]{0.45\textwidth}
\vspace{0pt}
\textbf{Check} each value in the \emph{original} equation.

\vspace{2pt}
\begin{tabular}{|l|}
\hline
$x=8$: \ $\ln(8^2-16)=\ln(6\cdot8)$ \\
\qquad\quad $\ln48=\ln48$ \ \checkmark \\ \hline
$x=-2$: \ $\ln((-2)^2-16)=\ln(6\cdot(-2))$ \\
\qquad\quad $\ln(-12)=\ln(-12)$ \ $\times$ \\ \hline
\end{tabular}

\vspace{6pt}
\textbf{A log with a coefficient?} Move the coefficient inside first, $k\log x=\log x^k$ for $x>0$ (\textbf{Power Property}), then set the arguments equal.
\end{minipage}

\vspace{4pt}
Logarithms are not defined for zero or negative arguments, so only $x=8$ is a solution. The value $x=-2$ is an \textbf{extraneous solution}: it solves the quadratic $x^2-6x-16=0$ but not the original log equation.

\textbf{Say it:} $x=-2$ is an extraneous solution because the argument $6x$ would be \blank[1.5cm], and \dots
\lines[2]

\wspart{Try It 5}{Solve each equation}
Hint: set the arguments equal, solve, and check every value in the original equation.

\begin{minipage}[t]{0.48\textwidth}
\textbf{a.} $\log_5(x^2-45)=\log_5(4x)$
\lines[5]
Solution: \blank[2.2cm] \quad Extraneous: \blank[1.8cm]
\end{minipage}\hfill
\begin{minipage}[t]{0.48\textwidth}
\textbf{b.} $\ln(-4x-1)=\ln(4x^2)$
\lines[5]
Solution: \blank[2.2cm] \quad Extraneous: \blank[1.8cm]
\end{minipage}

Is a negative solution always extraneous? Check both \emph{arguments}, not just $x$.
\lines[2]

\newpage

\wspart{IXL}{Practice for the quiz}
Do the starred skills first. \textbf{SmartScore 70 or above on each starred ($\bigstar$) skill replaces your quiz score.}

\vspace{2pt}
\begin{tabularx}{\textwidth}{l Y Y}
\toprule
\textbf{Code} & \textbf{Skill} & \textbf{Feeds} \\
\midrule
\ixlrow{YQY $\bigstar$}{Solve exponential equations by rewriting the base}{Quiz 2 and 3 (Example 1)}
\ixlrow{BXU $\bigstar$}{Solve logarithmic equations I}{Quiz 4 (Example 5) and Exit question b}
\ixlrow{RLX $\bigstar$}{Solve logarithmic equations II}{Quiz 4: a log with a coefficient in front}
\ixlrow{V2J}{Power rule}{Warm-up for quiz 3: the Power of a Power line}
\ixlrow{7T3}{Power property of logarithms}{Warm-up for quiz 4: $k\log x=\log x^k$}
\ixlrow{9F2}{Solve exponential equations using common logarithms}{Optional: quiz 1 (full Savvas quiz only)}
\bottomrule
\end{tabularx}

\wspart{Exit question}{Practice 43 (Higher Order Thinking)}
A biologist is using the logarithmic model $n=k\log(A)$ to determine the number of a species $n$ that can live on a land mass of area $A$. The constant $k$ varies according to the species.

\begin{minipage}[t]{0.50\textwidth}
\vspace{0pt}
\begin{tikzpicture}
\begin{axis}[sheetgrid, xmin=0, xmax=1350, ymin=0, ymax=1750,
  xtick={0,100,...,1300}, ytick={0,100,...,1700},
  xticklabels={0,,,,,500,,,,,1000,,,}, yticklabels={0,,,,,500,,,,,1000,,,,,1500,,},
  x=0.0053cm, y=0.0034cm, grid style={gray!35},
  xlabel={Acres}, ylabel={Species},
  xlabel style={at={(axis description cs:0.5,-0.07)}, anchor=north, font=\scriptsize},
  ylabel style={at={(axis description cs:-0.13,0.5)}, rotate=90, anchor=south, font=\scriptsize}]
\addplot[very thick, red, domain=0.05:3.114, samples=120, ->] ({10^x},{500*x});
\node[closed] at (axis cs:100,1000) {};
\node[closed] at (axis cs:1000,1500) {};
\node[font=\scriptsize, fill=white, inner sep=1pt, anchor=north west] at (axis cs:130,960) {$(100,1000)$};
\node[font=\scriptsize, fill=white, inner sep=1pt, anchor=north] at (axis cs:1000,1430) {$(1000,1500)$};
\end{axis}
\end{tikzpicture}
\end{minipage}\hfill
\begin{minipage}[t]{0.47\textwidth}
\vspace{0pt}
\textbf{a.} Use the graph to determine the constant $k$ for the species that the scientist is studying.
\lines[3]
\vspace{4pt}$k=\blank[2cm]$\par\vspace{12pt}

\textbf{b.} Determine the land mass in acres that is needed to support 3,000 of the species.
\lines[3]
$A=\blank[2.5cm]$ acres
\end{minipage}

\vspace{4pt}
\textbf{Before the quiz, finish these sentences.}

To solve $25^{3x}=125^{x+2}$, I first write both sides with the common base \blank[1.2cm], and then I set the \blank[2.2cm] equal.

To solve $\log_5(x^2-45)=\log_5(4x)$, I set the \blank[2.2cm] equal, and I check each answer because a value that makes an argument zero or negative gives an \blank[3.5cm].

\end{document}
